\documentclass{article}

\usepackage{graphicx}
\usepackage{amsmath, amssymb}

\title{Lecture 5 Scribe}
\author{Virja Shah}
\date{February 2026}

\begin{document}

\maketitle

\section*{Bayes' Theorem, Random Variables, and Probability Mass Function (PMF)}

\section{Bayes' Theorem}

\subsection{Weighted Average of Conditional Probabilities}

\subsubsection*{Step 1: Event Decomposition (Core Idea)}
Let $A$ and $B$ be events. We may express event $A$ as:
\[
A = AB \cup AB^c
\]

Reasoning (as stated in the slides):\\
For an outcome to be in $A$, it must either:
\begin{itemize}
\item be in both $A$ and $B$, or
\item be in $A$ but not in $B$.
\end{itemize}

\subsubsection*{Step 2: Mutual Exclusiveness}
The events $AB$ and $AB^c$ are mutually exclusive (cannot occur together).

\subsubsection*{Step 3: Apply Axiom 3 (Additivity for Disjoint Events)}
By Axiom 3:
\[
Pr(A)=Pr(AB)+Pr(AB^c)
\]

\subsubsection*{Step 4: Convert Intersections into Conditional Probability Form}
Using:
\[
Pr(AB)=Pr(A\mid B)Pr(B)
\]
and similarly,
\[
Pr(AB^c)=Pr(A\mid B^c)Pr(B^c)
\]

So:
\[
Pr(A)=Pr(A\mid B)Pr(B)+Pr(A\mid B^c)Pr(B^c)
\]

Since:
\[
Pr(B^c)=1-Pr(B)
\]

Final form:
\[
Pr(A)=Pr(A\mid B)Pr(B)+Pr(A\mid B^c)\big[1-Pr(B)\big]
\]

Interpretation (Exactly from slides):\\
The probability of event $A$ is a weighted average of the conditional probabilities,
with weights given by the probability of the event on which it is conditioned has of occurring.

\subsection{Learning by Example --- Example \#3.1 (Part 1/2)}

\subsubsection*{Problem Statement}
An insurance company divides people into:
\begin{itemize}
\item accident prone
\item not accident prone
\end{itemize}

Given:
\[
Pr(\text{accident in 1 year}\mid \text{accident prone})=0.4
\]
\[
Pr(\text{accident in 1 year}\mid \text{not accident prone})=0.2
\]
\[
Pr(\text{accident prone})=0.3
\]

Find: Probability a new policyholder has an accident within a year.

\subsubsection*{Step-by-step Solution (Exactly as in slides)}
Let:
\begin{itemize}
\item $A_1$ = event that policyholder has an accident within a year
\item $A$ = event that policyholder is accident prone
\end{itemize}

Then:
\[
Pr(A_1)=Pr(A_1\mid A)Pr(A)+Pr(A_1\mid A^c)Pr(A^c)
\]

Substitute values:
\[
Pr(A_1)=(0.4)(0.3)+(0.2)(0.7)
\]
\[
Pr(A_1)=0.12+0.14=0.26
\]

Final Answer:
\[
Pr(A_1)=0.26
\]

\subsection{Bayes Formula --- Example \#3.1 (Part 2/2)}

\subsubsection*{Problem Statement (Twist)}
Suppose the new policyholder has an accident within a year. Find:
\[
Pr(A\mid A_1)
\]

\subsubsection*{Step-by-step Solution (Exactly as in slides)}
\[
Pr(A\mid A_1)=\frac{Pr(AA_1)}{Pr(A_1)}
\]

Now:
\[
Pr(AA_1)=Pr(A)Pr(A_1\mid A)
\]

So:
\[
Pr(A\mid A_1)=\frac{Pr(A)Pr(A_1\mid A)}{Pr(A_1)}
\]

Substitute values:
\[
Pr(A\mid A_1)=\frac{(0.3)(0.4)}{0.26}
=\frac{0.12}{0.26}
=\frac{6}{13}
\]

Final Answer:
\[
Pr(A\mid A_1)=\frac{6}{13}
\]

\subsection{Law of Total Probability (Formula 3.4)}

\subsubsection*{Assumptions (Must be explicitly stated)}
Suppose $B_1,B_2,\dots,B_n$ are mutually exclusive events.

Their union forms the whole event space $B$:
\[
\bigcup_{i=1}^{n}B_i=B
\]

Meaning: Exactly one of the events $B_1,B_2,\dots,B_n$ must occur.

\subsubsection*{Derivation (As shown in slides)}
Write:
\[
A=\bigcup_{i=1}^{n}AB_i
\]

Since the events $AB_i$ are mutually exclusive:
\[
Pr(A)=\sum_{i=1}^{n}Pr(AB_i)
\]

Convert using conditional probability:
\[
Pr(AB_i)=Pr(A\mid B_i)Pr(B_i)
\]

Thus:
\[
Pr(A)=\sum_{i=1}^{n}Pr(A\mid B_i)Pr(B_i)
\]

This is the Law of Total Probability (Formula 3.4).

\subsection{Bayes Formula (Proposition 3.1)}

\subsubsection*{Key Identity Used}
The slides state:
\[
Pr(AB_i)=Pr(B_i\mid A)Pr(A)
\]

\subsubsection*{Bayes Formula (Final Result from slides)}
\[
Pr(B_i\mid A)=
\frac{Pr(A\mid B_i)Pr(B_i)}
{\sum_{j=1}^{n}Pr(A\mid B_j)Pr(B_j)}
\]

This is known as the Bayes Formula (Proposition 3.1).

\subsubsection*{Definitions (Explicit from slides)}
\[
Pr(B_i) = \text{apriori probability}
\]
probabilities that are formed from self-evident or presupposed models.

\[
Pr(B_i\mid A) = \text{posteriori probability}
\]
probabilities that are derived or calculated after observing certain events of event $B_i$ given $A$.

\subsection{Learning by Example --- Example \#3.2 (Cards)}

\subsubsection*{Problem Statement}
There are 3 cards:
\begin{enumerate}
\item RR (both sides red)
\item BB (both sides black)
\item RB (one side red, one side black)
\end{enumerate}

Cards are mixed, one card selected and placed down.\\
If the upper side is red, find probability that the other side is black.

\subsubsection*{Step-by-step Solution (Bayes Method from slides)}
Let:
\begin{itemize}
\item $RR,BB,RB$ denote events that chosen card is those types
\item $R$ = event that the upturned side is red
\end{itemize}

We want:
\[
Pr(RB\mid R)
\]

Bayes setup:
\[
Pr(RB\mid R)=\frac{Pr(RB\cap R)}{Pr(R)}
\]

Convert numerator:
\[
Pr(RB\cap R)=Pr(R\mid RB)Pr(RB)
\]

Denominator expands using total probability:
\[
Pr(R)=Pr(R\mid RR)Pr(RR)+Pr(R\mid RB)Pr(RB)+Pr(R\mid BB)Pr(BB)
\]

So:
\[
Pr(RB\mid R)=
\frac{Pr(R\mid RB)Pr(RB)}
{Pr(R\mid RR)Pr(RR)+Pr(R\mid RB)Pr(RB)+Pr(R\mid BB)Pr(BB)}
\]

Substitute values:
\[
Pr(RB)=Pr(RR)=Pr(BB)=\frac{1}{3}
\]
\[
Pr(R\mid RR)=1,\quad Pr(R\mid RB)=\frac{1}{2},\quad Pr(R\mid BB)=0
\]

Thus:
\[
Pr(RB\mid R)=
\frac{(1/2)(1/3)}
{(1)(1/3)+(1/2)(1/3)+(0)(1/3)}
\]

Compute:
\[
\text{Numerator}=\frac{1}{6}
\]
\[
\text{Denominator}=\frac{1}{3}+\frac{1}{6}+0=\frac{1}{2}
\]

Thus:
\[
Pr(RB\mid R)=\frac{1/6}{1/2}=\frac{1}{3}
\]

Final Answer:
\[
Pr(RB\mid R)=\frac{1}{3}
\]

\subsubsection*{Slide Explanation of Common Mistake}
Many guess $\frac{1}{2}$ incorrectly by assuming two equally likely cases (RR or RB).

The slides explain the correct reasoning by listing 6 equally likely sides:
\[
R_1,R_2,B_1,B_2,R_3,B_3
\]

Given a red side, only outcomes are:
\[
R_1,R_2,R_3
\]

Only $R_3$ corresponds to RB card having black on the other side.\\
So probability $=\frac{1}{3}$.

\section{Random Variables}

\subsection{Motivation and Concept}

Key Motivation (From slides):\\
When an experiment is performed, we are frequently interested mainly in some function of the outcome rather than the actual outcome itself.

Examples from slides:
\begin{itemize}
\item Dice tossing: focus on sum of dice rather than individual values.
\item Coin flipping: interest may lie in total number of heads rather than exact sequence.
\end{itemize}

\subsection{Definition (From slide)}
These real-valued functions defined on the sample space are called Random Variables.

Properties (explicit from slide):
\begin{itemize}
\item Values are determined by outcomes of an experiment.
\item Probabilities are assigned to possible values of Random Variables.
\end{itemize}

\subsection{Formal Definition (From slide)}
A random variable $X$ on a sample space $\Omega$ is a function:
\[
X:\Omega\to \mathbb{R}
\]

It assigns to each sample point $\omega\in\Omega$ a real number:
\[
X(\omega)
\]

\subsection{Discrete Restriction (Assumption stated in slides)}
Until further notice, attention is restricted to discrete random variables, meaning:
\begin{itemize}
\item values are in a range that is finite or countably infinite
\item even though $X$ maps into $\mathbb{R}$, the set of values $X$ takes is a discrete subset of $\mathbb{R}$.
\end{itemize}

\section{Distribution of a Random Variable}

Probability space defines:
\begin{itemize}
\item sample space $\Omega$
\item probability of each sample point
\end{itemize}

Analogously, random variable defines:
\begin{itemize}
\item set of values it can take
\item probabilities with which it takes those values
\end{itemize}

\subsection{Event Form}
Let $a$ be a number in the range of $X$. Then:
\[
\{\omega\in\Omega:X(\omega)=a\}
\]

This is an event in the sample space.

Abbreviated as:
\[
X=a
\]

Then:
\[
Pr[X=a]
\]

The collection of these probabilities for all values of $a$ is called the distribution of $X$.

\subsection{Visualization (Bar Diagram)}
The distribution can be visualized as a bar diagram:
\begin{itemize}
\item x-axis: values random variable can take
\item bar height at value $a$: $Pr[X=a]$
\item each probability computed by probability of corresponding event in sample space
\end{itemize}

\subsection{Random Variable Example (Example-1)}

Experiment: toss 3 fair coins.\\
Let $Y$ = number of heads.

Possible values:
\[
Y\in\{0,1,2,3\}
\]

Step-by-step Probabilities (From slides):
\[
P\{Y=0\}=P\{(t,t,t)\}=\frac{1}{8}
\]
\[
P\{Y=1\}=P\{(t,t,h),(t,h,t),(h,t,t)\}=\frac{3}{8}
\]
\[
P\{Y=2\}=P\{(t,h,h),(h,t,h),(h,h,t)\}=\frac{3}{8}
\]
\[
P\{Y=3\}=P\{(h,h,h)\}=\frac{1}{8}
\]

\subsection{Total Probability Check (From slides)}
Since $Y$ must take one of the values 0 through 3:
\[
1=P\left(\bigcup_{i=0}^{3}\{Y=i\}\right)
=\sum_{i=0}^{3}P(Y=i)
\]

\section{Probability Mass Function (PMF)}

\subsection{Discrete Random Variable Definition}
A random variable that can take on at most a countable number of possible values is said to be discrete.

\subsection{PMF Definition (Formal)}
Let $X$ be a discrete random variable with range:
\[
R_X=x_1,x_2,x_3,\dots
\]

The function:
\[
p(x_k)=Pr(X=x_k)
\]
is called the Probability Mass Function (PMF) of $X$.

\subsection{PMF Property (Normalization)}
Since $X$ must take on one of the values $x_k$, we have:
\[
\sum_k p(x_k)=1
\]

\subsection{PMF Example (Given in slides)}

A random variable $X$ has PMF:
\[
p(i)=c\frac{\lambda^i}{i!},\quad i=0,1,2,\dots
\]
where $\lambda$ is positive.

Find:
\[
P\{X=0\},\quad P\{X>2\}
\]

\subsection{Step 1: Use the PMF Normalization Condition}
Since:
\[
\sum_{i=0}^{\infty}p(i)=1
\]

Substitute:
\[
c\sum_{i=0}^{\infty}\frac{\lambda^i}{i!}=1
\]

\subsection{Step 2: Use Series Identity from slide}
The slide states:
\[
e^\lambda=\sum_{i=0}^{\infty}\frac{\lambda^i}{i!}
\]

So:
\[
ce^\lambda=1
\]

Thus:
\[
c=e^{-\lambda}
\]

\subsection{Step 3: Compute $P\{X=0\}$}
\[
P\{X=0\}=p(0)=c\frac{\lambda^0}{0!}
\]

Substitute:
\[
c=e^{-\lambda},\quad \lambda^0=1,\quad 0!=1
\]

Final Answer:
\[
P\{X=0\}=e^{-\lambda}
\]

\subsection{Step 4: Compute $P\{X>2\}$}

Slides use complement:
\[
P\{X>2\}=1-P\{X\le 2\}
\]

Then:
\[
P\{X>2\}=1-P\{X=0\}-P\{X=1\}-P\{X=2\}
\]

\section*{Final Exam Revision Checklist (From Lecture 5)}

Bayes \& Total Probability:
\[
A=AB\cup AB^c
\]
\[
Pr(A)=Pr(A\mid B)Pr(B)+Pr(A\mid B^c)\big[1-Pr(B)\big]
\]
\[
Pr(A)=\sum_{i=1}^{n}Pr(A\mid B_i)Pr(B_i)
\]
\[
Pr(B_i\mid A)=
\frac{Pr(A\mid B_i)Pr(B_i)}
{\sum_{j=1}^{n}Pr(A\mid B_j)Pr(B_j)}
\]

Random Variables:
\[
X:\Omega\to \mathbb{R}
\]
Distribution:
\[
Pr[X=a]
\]

PMF:
\[
p(x_k)=Pr(X=x_k)
\]
\[
\sum_k p(x_k)=1
\]

\end{document}
